\section{Supervised Contrastive Learning for Enhanced Representation Learning}

\subsection{Introduction}

Supervised Contrastive Learning (SCL), as introduced by Khosla et al. \cite{khosla2020supervised}, is a powerful machine learning technique within the realm of deep learning, specifically designed to enhance the process of representation learning for supervised classification tasks. Building upon the foundations of Contrastive Learning (CL) \cite{hadsell2006dimensionality, chen2020simple}, SCL leverages the availability of class labels to create more discriminative feature embeddings. By explicitly contrasting data points based on their assigned labels, SCL enables models to learn feature spaces where instances of the same class are tightly clustered together, while instances from different classes are well-separated.

The primary objective of Supervised Contrastive Learning is to learn feature representations such that data samples belonging to the same class (sharing the same label) exhibit high similarity and thus are located close to each other in the embedding space. Conversely, data samples from different classes (having different labels) are pushed apart, resulting in low similarity and greater distance in the embedding space. This learning objective leads to the development of a robust and highly discriminative feature space, where the data points of each class form compact clusters that are well-isolated from other classes, ultimately improving the performance of downstream classification tasks.

\subsection{Operational Principles}

To achieve the goal of learning discriminative feature representations, Supervised Contrastive Learning employs a contrastive loss function that is adapted to utilize the available label information. This loss function compares data points within the embedding space and optimizes the embeddings such that points with the same label are drawn closer together, while points with different labels are pushed farther apart. The Supervised Contrastive Loss can be mathematically expressed as follows:

$$
\mathcal{L}_{SCL} = \sum_{i \in N} \frac{-1}{|P(i)|} \sum_{p \in P(i)} \log \frac{\exp(\text{sim}(z_i, z_p) / \tau)}{\sum_{a \in A(i)} \exp(\text{sim}(z_i, z_a) / \tau) + \exp(\text{sim}(z_i, z_i) / \tau)}
$$

where:
\begin{itemize}
    \item $N$ is the total number of samples in the batch.
    \item $P(i)$ is the set of indices of all data points in the batch that have the same label as the $i$-th data point, excluding $i$ itself. Thus, $|P(i)|$ is the number of positive pairs for the $i$-th sample within the batch.
    \item $A(i)$ is the set of indices of all data points in the batch that have a different label from the $i$-th data point. These represent the negative samples for the $i$-th data point within the batch.
    \item $\text{sim}(z_i, z_j)$ is a similarity function that measures the similarity between the feature vectors $z_i$ and $z_j$ of samples $i$ and $j$. Cosine similarity is a commonly used similarity metric:
    $$
    \text{sim}(z_i, z_j) = \frac{z_i \cdot z_j}{\|z_i\| \cdot \|z_j\|}
    $$
    where $z_i$ and $z_j$ are typically L2 normalized.
    \item $\tau$ is the temperature parameter, a positive scalar that controls the sharpness of the contrastive distribution. A lower temperature makes the loss more sensitive to hard negative samples.
    \item $z_i$ is the feature vector (embedding) of the $i$-th data sample, generated by a deep neural network.
\end{itemize}

The Supervised Contrastive Loss operates on the principle of encouraging clustering of samples with the same label while promoting separation between samples with different labels within the embedding space. For each anchor point $i$ in the batch, the loss considers all other points with the same label as positive samples and all points with different labels as negative samples. The objective is to maximize the similarity between the anchor and its positive samples (numerator of the log term) while minimizing the similarity between the anchor and its negative samples (denominator of the log term). The normalization by $|P(i)|$ averages the loss over all positive pairs for each anchor. The inclusion of $\exp(\text{sim}(z_i, z_i) / \tau)$ in the denominator accounts for the similarity of the sample with itself, although its contribution is constant and primarily ensures numerical stability.

\subsection{Comparison with Unsupervised Contrastive Learning}

Supervised Contrastive Learning builds upon the principles of unsupervised Contrastive Learning (UCL) methods like SimCLR \cite{chen2020simple} and MoCo \cite{he2020momentum}, but with a crucial distinction: the utilization of label information. Here's a comparison of the two paradigms:

\subsubsection{Unsupervised Contrastive Learning}

\textbf{Data Augmentation as Supervision:} UCL relies on generating multiple augmented views of the same unlabeled data instance to create positive pairs. The underlying assumption is that different augmented views of the same instance should have similar representations. Negative pairs are formed by augmented views of different data instances within the batch.
\textbf{Loss Function:} The loss function in UCL aims to maximize the similarity between different augmented views of the same instance and minimize the similarity between augmented views of different instances. A common form of the loss for an anchor $i$ and its positive counterpart $j$ is:
    $$
    \mathcal{L}_{UCL}^{(i)} = - \log \frac{\exp(\text{sim}(z_i, z_j) / \tau)}{\sum_{k \neq i} \exp(\text{sim}(z_i, z_k) / \tau)}
    $$
    where $z_k$ represents the embeddings of all other augmented views in the batch (including other positive pairs).
\textbf{Objective:} The primary objective is to learn general-purpose representations that capture the inherent structure and variations within the unlabeled data. These representations can then be used for various downstream tasks, often requiring fine-tuning with labeled data.

\subsubsection{Supervised Contrastive Learning}

\textbf{Label-Based Positive Pairs:} SCL directly utilizes class labels to define positive pairs. Instead of relying solely on augmentations of the same instance, SCL considers all instances within the same batch that share the same class label as positive pairs for a given anchor instance.

\textbf{More Informative Negative Samples:} The negative samples in SCL are explicitly defined as instances belonging to different classes within the batch. This provides more direct and informative negative examples for the model to discriminate against.

\textbf{Loss Function:} As shown in the previous section, the Supervised Contrastive Loss explicitly sums over all positive pairs for each anchor and contrasts them against all negative samples in the batch.

\textbf{Objective:} The primary objective of SCL is to learn representations that are specifically tailored for classification. By leveraging label information during training, SCL aims to create a feature space where class separability is maximized, leading to improved performance on supervised classification tasks, often with less reliance on extensive data augmentation compared to UCL for downstream tasks.

\textbf{Key Differences and Advantages of SCL:}

\textbf{Stronger Supervision:} SCL benefits from the explicit supervision provided by class labels, leading to more discriminative representations compared to UCL, which relies on the inductive bias that different augmentations of the same instance should be similar.
\textbf{Improved Class Separability:} By explicitly pulling together instances of the same class and pushing apart instances of different classes, SCL directly optimizes for class separability in the embedding space.
\textbf{Potentially Less Reliance on Extensive Augmentations:} While augmentations can still be beneficial in SCL to improve robustness, the strong signal from the labels can reduce the need for very diverse and aggressive augmentations that are often crucial for the success of UCL.
\textbf{Task-Specific Representations:} The representations learned by SCL are more directly aligned with the classification task for which the labels are available, often leading to better performance on these tasks compared to the more general-purpose representations learned by UCL.

In conclusion, Supervised Contrastive Learning represents a significant advancement over unsupervised contrastive learning for supervised tasks by effectively incorporating label information into the contrastive learning framework. This leads to the learning of more discriminative and task-specific feature representations, ultimately enhancing the performance of classification models.

